% !PRINTED 642
\weberpage{642}

4. \gesperrt{Jedes ganze rationale Funktional läßt sich in Primfaktoren zerlegen und zwar, von Einheitsfaktoren abgesehen, nur auf eine Weise.}

\gesperrt{Ist ein Produkt von ganzen Funktionalen in $\overline{Z}$ durch ein Primfunktional teilbar, so ist wenigstens einer der Faktoren durch dieses Primfunktional teilbar.}

\begin{center}
§ 178.\\
\gesperrt{Funktionale des Körpers $\overline{\Omega}$.}
\end{center}

Ein Funktional des Körpers $\overline{\Omega}$ ist ein Ausdruck von der Form:
\weq{1}{\overline{\omega} = \overline{x}_0 + \overline{x}_1\theta + \cdots + \overline{x}_{n-1}\theta^{n-1},}
wo $\theta$ die den Körper $\Omega$ erzeugende algebraische Funktion ist, und $\overline{x}_0, \overline{x}_1, \ldots, \overline{x}_{n-1}$ \gesperrt{rationale Funktionale} sind. $\overline{\omega}$ ist nur dann $= 0$, wenn alle Koeffizienten $\overline{x}_0, \overline{x}_1, \ldots, \overline{x}_{n-1}$ verschwinden.

Jedes Funktional, das in bezug auf $\theta$ von höherem als dem $(n-1)$ten Grad ist, kann durch den Rest der Division durch $f(\theta)$ ersetzt, also auf den $(n-1)$ten Grad reduziert werden, und der Quotient zweier Ausdrücke (1), in der der Nenner nicht $=0$ ist, kann wie in § 171, 1. auf die Form (1) gebracht werden.
% [sic] printed "in der der Nenner", the relative pronoun should be "dem" (antecedent "der Quotient")

Nimmt man eine beliebige Basis
\[
\omega_1, \omega_2, \ldots, \omega_n
\]
des Körpers $\Omega$, so kann $\overline{\omega}$ auch in die Form
\weq{2}{\overline{\omega} = \overline{x}_1\omega_1 + \overline{x}_2\omega_2 + \cdots + \overline{x}_n\omega_n}
gesetzt werden. Wendet man dies auf die Produkte $\overline{\omega}\omega_i$ an und setzt
\weq{3}{\overline{\omega}\omega_i = \overline{x}_{i,1}\omega_1 + \overline{x}_{i,2}\omega_2 + \cdots + \overline{x}_{i,n}\omega_n,}
worin die $\overline{x}_{i,n}$ rationale Funktionale sind, so ergibt sich für $\overline{\omega}$ eine Gleichung:
\weq{4}{\left|\begin{array}{lll}
\overline{x}_{1,1} - \overline{\omega}, & \overline{x}_{1,2}, & \ldots, \overline{x}_{1,n} \\
\overline{x}_{2,1}, & \overline{x}_{2,2} - \overline{\omega}, & \ldots, \overline{x}_{2,n} \\
\cdot & \cdot & \cdots \\
\overline{x}_{n,1}, & \overline{x}_{n,2}, & \ldots, \overline{x}_{n,n} - \overline{\omega}
\end{array}\right| = 0,}
also:

5. \gesperrt{Jedes Funktional in $\overline{\Omega}$ genügt einer Gleichung}
\weq{5}{\varphi(\overline{\omega}) = 0,}
\gesperrt{worin}
\weq{6}{\varphi(t) = t^n + \overline{A}_1 t_2^{n-1} + \overline{A}_2 t^{n-2}, \cdots + \overline{A}_n}
% [sic] printed "\overline{A}_1 t_2^{n-1}" and "t^{n-2}, \cdots", misprints for "\overline{A}_1 t^{n-1}" (stray subscript 2 on t) and "t^{n-2} + \cdots" (comma for plus, cf. (1)-(3))
%CONT>>
